\documentclass[a4paper,11pt]{article}
\usepackage[utf8]{inputenc}
\usepackage[T1]{fontenc}
\usepackage[french]{babel}
\usepackage[left=2cm,right=2cm,top=2cm,bottom=2cm]{geometry}
\usepackage{xcolor}
\usepackage{hyperref}
\usepackage{fontawesome5}
\usepackage{parskip}

% --- Couleurs ---
\definecolor{primary}{RGB}{0, 82, 147}
\hypersetup{colorlinks=true, urlcolor=primary}

\begin{document}

% --- EXPÉDITEUR ---
\noindent
\begin{minipage}[t]{0.5\textwidth}
    \textbf{Loïc Aron MBASSI EWOLO}\\
    Élève-ingénieur ENSEA / Polytechnique Yaoundé\\[3pt]
    \faMapMarker\ Cergy, France\\
    \faPhone\ (+33) 7 59 52 33 98\\
    \faEnvelope\ \href{mailto:loicaron.mbassiewolo@ensea.fr}{loicaron.mbassiewolo@ensea.fr}
\end{minipage}
\hfill
% --- DATE ---
\begin{minipage}[t]{0.4\textwidth}
    \raggedleft
    Cergy, le 5 février 2026
\end{minipage}

\vspace{1.2cm}

% --- DESTINATAIRE ---
\hfill
\begin{minipage}[t]{0.5\textwidth}
    \raggedleft
    \textbf{STMicroelectronics}\\
    Service Recrutement\\
    Grenoble, France
\end{minipage}

\vspace{1cm}

\noindent\textbf{Objet :} Candidature au stage Ingénieur Développement Python (Réf. 9234)

\vspace{0.8cm}

Madame, Monsieur,

L'an dernier, j'ai passé des semaines à modéliser un drone sous MATLAB/Simulink pour un projet de détection de défauts. Le plus frustrant n'était pas la modélisation en soi, mais les heures perdues à relancer manuellement les simulations avec différents paramètres. C'est en lisant votre offre sur l'automatisation des simulations Monte Carlo que j'ai compris que ce problème était universel dans l'industrie.

À Polytechnique Yaoundé, ma formation en mathématiques appliquées m'a donné les bases : algèbre linéaire, analyse numérique, probabilités. J'ai appris à coder proprement en Python, pas juste à faire tourner des scripts. Mon outil open source de génération de code (Laravel API Generator) m'a obligé à structurer du code maintenable, à parser des fichiers XML avec des calculs géométriques, et à documenter pour que d'autres puissent l'utiliser. Ce n'est pas de la simulation électromagnétique, mais la logique de workflow automatisé est la même.

Ce qui m'attire dans ce stage, c'est le côté concret. Vous partez d'un problème réel (la variabilité de fabrication) et vous construisez un outil pour y répondre. J'ai fait ça à plus petite échelle avec mon projet d'analyse ECG : nettoyer les données, extraire des features, évaluer les performances du modèle. Je n'ai jamais utilisé COMSOL, mais j'apprends vite quand le problème m'intéresse.

Je suis disponible à partir d'avril 2026 pour 5 mois. Si mon profil correspond à ce que vous cherchez, je serai ravi d'en discuter.

\vspace{0.6cm}

Cordialement,

\vspace{0.4cm}

\textbf{Loïc Aron MBASSI EWOLO}

\end{document}
